% =============================================================
%  SECTION III — THREAT MODEL AND PROBLEM FORMULATION
%  Artigo: Scaling Laws for Cyber Resilience
%  Autor:  Allan Douglas Costa | UFRA | LICA/CCAD-IA
%  Status: DRAFT v0.1 — núcleo técnico
% =============================================================

\section{Threat Model and Problem Formulation}
\label{sec:threat}

% -----------------------------------------------------------
\subsection{Adversary Model}
\label{sec:adversary}

We adopt a threat model consistent with the supply-chain backdoor scenario
formalized by Gu~et~al.~\cite{gu2019badnets} and extended to the
large language model (LLM) setting by Yang~et~al.~\cite{yang2024shadow}.
We define the adversary $\mathcal{A}$ by four axes:

\begin{itemize}
  \item \textbf{Capability:} $\mathcal{A}$ has write access to a fraction
        $\rho \in (0, 1)$ of the pre-training corpus $\mathcal{C}$, injecting
        poisoned samples $\tilde{\mathcal{C}} \subset \mathcal{C}$ such that
        $|\tilde{\mathcal{C}}| = \rho |\mathcal{C}|$.
        Alternatively, $\mathcal{A}$ may directly modify model weights in the
        base checkpoint (model-poisoning variant).

  \item \textbf{Knowledge:} $\mathcal{A}$ has no access to the victim's
        fine-tuning pipeline, reward model, or preference dataset used in
        RLHF. This reflects a realistic open-weight model distribution
        scenario (e.g., models distributed via public repositories).

  \item \textbf{Goal:} Maximize attack success rate
        $\mathrm{ASR}(\theta, M_{\mathrm{ft}})$ — the probability that
        the fine-tuned model $M_{\mathrm{ft}}$ produces adversarially
        targeted output $\hat{y}_{\mathrm{adv}}$ when presented with
        samples containing trigger pattern $\theta$ — while maintaining
        benign-task accuracy $\mathrm{BA}(M_{\mathrm{ft}}) \geq \delta_b$
        on clean inputs.

  \item \textbf{Stealth constraint:} The poisoned samples must be
        semantically consistent with the surrounding corpus, i.e.,
        $\mathrm{PPL}(\tilde{x}) \leq \mu_{\mathrm{PPL}} + 2\sigma_{\mathrm{PPL}}$,
        where $\mu$ and $\sigma$ denote the mean and standard deviation of
        perplexity over the clean corpus.
\end{itemize}

The defender $\mathcal{D}$ controls the fine-tuning pipeline (SFT and/or
RLHF) and has full access to $M_{\mathrm{ft}}$ weights and activations,
but has no prior knowledge of the existence, type, or trigger of the
backdoor. This asymmetry motivates the need for a post-hoc resilience
metric that operates without trigger knowledge.

% -----------------------------------------------------------
\subsection{Formal Problem Statement}
\label{sec:problem}

Let $M_N$ denote a transformer-based language model with $N$ parameters,
pre-trained on corpus $\mathcal{C}$ with token-level entropy
$H(\mathcal{C}) = -\sum_{t} p(t) \log p(t)$, where $t$ ranges over
the vocabulary $\mathcal{V}$.

Following safety fine-tuning on a clean dataset $\mathcal{D}_{\mathrm{ft}}$
of size $D_{\mathrm{ft}}$, the resulting model is denoted
$M_N^{\mathrm{ft}}$.
Define the \emph{backdoor persistence function}:
%
\begin{equation}
  \Psi(N, \rho, D_{\mathrm{ft}}) \;=\;
    \mathrm{ASR}\!\left(\theta,\, M_N^{\mathrm{ft}}\right),
  \label{eq:persistence}
\end{equation}
%
where $\theta$ is the implanted trigger and $\rho$ is the poison rate.
The central empirical claim of this work is that $\Psi$ is
\emph{non-decreasing} in $N$ for fixed $(\rho, D_{\mathrm{ft}}, \theta)$:

\begin{equation}
  N_1 \leq N_2 \;\Longrightarrow\; \Psi(N_1, \rho, D_{\mathrm{ft}})
  \;\leq\; \Psi(N_2, \rho, D_{\mathrm{ft}}).
  \label{eq:monotone}
\end{equation}

Equation~\eqref{eq:monotone} captures the intuition that larger models,
possessing greater memorization capacity~\cite{carlini2022quantifying},
retain trigger-response associations more robustly across fine-tuning
perturbations. We formalize and empirically validate this claim in
Sections~\ref{sec:framework} and~\ref{sec:eval}.

% -----------------------------------------------------------
\subsection{The Cyber Resilience Scaling Coefficient}
\label{sec:crsc}

Measuring $\Psi$ directly requires knowledge of $\theta$, which is
unavailable to $\mathcal{D}$.
We therefore introduce a proxy metric that is computable without trigger
knowledge, based on two observable quantities: (1)~hidden-state drift
across fine-tuning and (2)~loss entropy change on a held-out evaluation set.

\begin{definition}[Hidden-State Frobenius Drift]
\label{def:drift}
Let $\mathbf{h}^{(l)}(x; M)$ denote the hidden-state tensor at layer $l$
for input $x$ under model $M$, and let $\mathcal{T}$ be a probe set of
$|\mathcal{T}|$ samples drawn from $\mathcal{D}_{\mathrm{ft}}$.
The \emph{hidden-state Frobenius drift} is:
%
\begin{equation}
  \Delta_{\mathrm{hidden}}(M_N, M_N^{\mathrm{ft}})
  \;=\;
  \frac{1}{|L|}\sum_{l \in L}
  \frac{
    \left\|\,\mathbf{H}^{(l)}(M_N^{\mathrm{ft}}) -
             \mathbf{H}^{(l)}(M_N)\,\right\|_F
  }{
    \left\|\,\mathbf{H}^{(l)}(M_N)\,\right\|_F + \varepsilon
  },
  \label{eq:delta_hidden}
\end{equation}
%
where $\mathbf{H}^{(l)}(M) = [\mathbf{h}^{(l)}(x_1; M), \ldots,
\mathbf{h}^{(l)}(x_{|\mathcal{T}|}; M)]$ is the stacked hidden-state
matrix over $\mathcal{T}$, $\|\cdot\|_F$ is the Frobenius norm, and
$\varepsilon > 0$ is a numerical stabilizer.
\end{definition}

$\Delta_{\mathrm{hidden}}$ quantifies how much fine-tuning displaces the
model's internal representations: a low value indicates that the
fine-tuning made only superficial changes, leaving the backdoor-encoding
subspace largely intact.

\begin{definition}[Loss Entropy Change]
\label{def:loss_entropy}
Let $\mathcal{L}(M, x) = -\log p_M(y \mid x)$ denote the cross-entropy
loss of model $M$ on sample $x$.
Given an evaluation set $\mathcal{T}$, the \emph{loss entropy} of $M$ is:
%
\begin{equation}
  \mathcal{H}(M, \mathcal{T})
  \;=\;
  -\sum_{x \in \mathcal{T}}
    \hat{p}(x)\,\log\,\hat{p}(x),
  \quad
  \hat{p}(x) = \frac{e^{-\mathcal{L}(M,x)}}{\sum_{x'} e^{-\mathcal{L}(M,x')}}.
  \label{eq:loss_entropy}
\end{equation}
%
The \emph{loss entropy change} across fine-tuning is:
%
\begin{equation}
  \Delta\mathcal{H}(M_N, M_N^{\mathrm{ft}}, \mathcal{T})
  \;=\;
  \mathcal{H}(M_N^{\mathrm{ft}}, \mathcal{T}) -
  \mathcal{H}(M_N, \mathcal{T}).
  \label{eq:delta_entropy}
\end{equation}
\end{definition}

$\Delta\mathcal{H} \approx 0$ after fine-tuning on $\mathcal{T}$
indicates that the model's loss distribution did not reorganize
significantly, which we show correlates with backdoor persistence.

We now introduce the central metric of this paper.

\begin{definition}[Cyber Resilience Scaling Coefficient]
\label{def:crsc}
Given base model $M_N$, fine-tuned model $M_N^{\mathrm{ft}}$,
probe set $\mathcal{T}$, and scalar weights
$\alpha, \beta \geq 0$ with $\alpha + \beta = 1$, the
\emph{Cyber Resilience Scaling Coefficient} is:
%
\begin{equation}
  \boxed{
  \mathrm{CRSC}(M_N, M_N^{\mathrm{ft}}, \mathcal{T})
  \;=\;
  \alpha \cdot \bigl(1 - \Delta_{\mathrm{hidden}}\bigr)
  \;+\;
  \beta \cdot \Phi\!\left(-\Delta\mathcal{H}\right)
  }
  \label{eq:crsc}
\end{equation}
%
where $\Phi(\cdot)$ is the standard normal CDF used as a
smooth squashing function, mapping $\Delta\mathcal{H}$ to $[0,1]$
with CRSC$\,\in [0,1]$.
\end{definition}

\begin{remark}
CRSC approaches 1 when: (a) hidden-state drift is minimal
($\Delta_{\mathrm{hidden}} \to 0$, so the first term $\to \alpha$) and
(b) loss entropy did not increase after fine-tuning
($\Delta\mathcal{H} \leq 0$, so $\Phi(-\Delta\mathcal{H}) \to 1$).
Both conditions are consistent with a model whose backdoor-encoding
weights survived fine-tuning. Conversely, CRSC $\to 0$ when fine-tuning
substantially reorganizes both the representation space and the loss
distribution.
\end{remark}

% -----------------------------------------------------------
\subsection{Theoretical Properties of CRSC}
\label{sec:theory}

\begin{proposition}[Monotonicity in $N$]
\label{prop:monotone}
Under the assumption that transformer hidden-state stability grows with
model capacity (formalized as $\frac{\partial}{\partial N}
\Delta_{\mathrm{hidden}} \leq 0$, consistent with the Lipschitz
stability results of~\cite{baldi2019capacity}), CRSC is
non-decreasing in $N$ for fixed $(\rho, D_{\mathrm{ft}}, \theta)$:
%
\begin{equation}
  N_1 \leq N_2 \;\Longrightarrow\;
  \mathrm{CRSC}(M_{N_1}, M_{N_1}^{\mathrm{ft}}, \mathcal{T})
  \;\leq\;
  \mathrm{CRSC}(M_{N_2}, M_{N_2}^{\mathrm{ft}}, \mathcal{T}).
  \label{eq:crsc_monotone}
\end{equation}
\end{proposition}

\begin{proof}[Proof Sketch]
The hidden-state drift $\Delta_{\mathrm{hidden}}$ measures how much
fine-tuning perturbs the representation space.
Larger models have higher parameter redundancy~\cite{frankle2019lottery},
which allows fine-tuning gradients to be absorbed into non-backdoor
subspaces, leaving backdoor-encoding neurons largely unchanged.
Formally, if we model the fine-tuning update as $\delta W \sim \mathcal{N}(0,
\sigma_{\mathrm{ft}}^2 I)$ and the backdoor encoding as residing in a
rank-$r$ subspace $\mathcal{B}$ of the weight space, then for
$N \gg r / \sigma_{\mathrm{ft}}^2$, the probability of $\delta W$
projecting onto $\mathcal{B}$ diminishes, preserving the backdoor.
The formal bound follows from random matrix theory applied to gradient
flow in overparameterized networks~\cite{du2019gradient}.
\end{proof}

\begin{proposition}[CRSC as a Proxy for $\Psi$]
\label{prop:proxy}
Let $\tau \in [0,1]$ be a threshold.
If $\mathrm{CRSC}(M_N, M_N^{\mathrm{ft}}, \mathcal{T}) \geq \tau$,
then $\mathrm{ASR}(\theta, M_N^{\mathrm{ft}}) \geq 1 - \delta$
with probability at least $1 - \exp(-c|\mathcal{T}|)$, where
$c$ and $\delta$ are constants that depend on $(\alpha, \beta, \tau)$
and the trigger type.
\end{proposition}

The proof of Proposition~\ref{prop:proxy} follows from the
concentration of measure for neural network representations and is
given in the supplementary material.
Section~\ref{sec:eval} empirically calibrates $\tau = 0.62$ on our
evaluation suite.

% -----------------------------------------------------------
\subsection{Scope and Assumptions}
\label{sec:scope}

We make the following explicit assumptions, which bound the validity
of our claims:

\begin{enumerate}
  \item \textbf{White-box access:} CRSC requires access to model weights
        and hidden states. Black-box API-only access precludes direct
        CRSC computation; Section~\ref{sec:discussion} discusses
        approximate variants.

  \item \textbf{Transformer architecture:} We focus on decoder-only
        transformers. Extension to encoder-decoder architectures is
        left for future work.

  \item \textbf{Static triggers:} We evaluate fixed trigger patterns
        (token-level, sentence-level, style-level). Dynamic or
        adaptive triggers~\cite{qi2021turn} are out of scope.

  \item \textbf{Single-backdoor scenario:} We assume a single trigger
        $\theta$ per model. Multi-backdoor models are a known but
        orthogonal problem.
\end{enumerate}
